\section{Evaluierung und Ergebnisse}

\subsection{Evaluierungsmetriken und Test-Setup}
Dieses Kapitel definiert die quantitativen Error-Metriken (MSE, MAE) zur Leistungsmessung und etabliert den Bewertungsmaßstab auf dem isolierten Test-Set.

\subsection{Ergebnisse des Behavior Cloning}

\subsubsection{Quantitative Loss-Analyse}
Es erfolgt die Auswertung und Interpretation der berechneten Metriken sowie der Train- vs. Validation-Loss-Kurven zum Nachweis verhinderter Überanpassung.

\subsubsection{Visuelle Verhaltensanalyse}
Dieser Abschnitt stellt die vorhergesagten Actions (Hebelstellungen) und die Ground Truth (Fahrereingaben) grafisch über die Zeit gegenüber.

\subsection{Vergleichende RL-Evaluierung (Warm Start vs. Cold Start)}
Es werden die Reward-Kurven eines Agenten mit und ohne Pre-Training in der Simulation gegenübergestellt und der zeitliche sowie qualitative Performance-Gewinn quantifiziert.

\subsection{Diskussion und Limitationen}
Dieses Kapitel interpretiert die erzielten Ergebnisse objektiv und ordnet verbleibende Fehlerquellen (Covariate Shift) sowie Limitationen im Datensatz kritisch ein.
